\section*{Aktivitas Pembelajaran dan Praktikum}
\addcontentsline{toc}{section}{Aktivitas Pembelajaran dan Praktikum}

\begin{table}[H]
\centering
\caption{Rencana aktivitas Bab I}
\label{tab:ch01-aktivitas}
\begin{tabularx}{\textwidth}{@{}llX@{}}
\toprule
\textbf{Sesi} & \textbf{Bentuk} & \textbf{Kegiatan} \\
\midrule
1 & Ceramah & Membaca peta buku dan RPS; diskusi capaian pembelajaran \\
2 & Latihan & Menyusun rencana belajar pribadi dan memeriksa prasyarat \\
3 & Praktikum 0 & Menyiapkan lingkungan kerja dan menjalankan program contoh \\
\bottomrule
\end{tabularx}
\end{table}

\subsection*{Praktikum 0: Persiapan Lingkungan}
\begin{enumerate}[label=\arabic*.,leftmargin=*]
  \item Pastikan \code{pdflatex} dan \code{bibtex} tersedia.
  \item Unduh atau salin direktori proyek \file{code/subsetc}.
  \item Baca \file{code/subsetc/README.md} dan spesifikasi bahasa.
  \item Jalankan program contoh \file{hello.sc} melalui alat yang
        disediakan dosen.
  \item Tuliskan satu paragraf refleksi tentang kesulitan yang dihadapi.
\end{enumerate}

Praktikum ini tidak dinilai secara sumatif, tetapi menjadi prasyarat
administratif untuk Praktikum 1.
